% ============================================================================
% Fichier  : partie4/P04_C16_forensique_mobile.tex
% Titre    : Forensique Mobile -- Android et iOS
% Auteur   : MINKA MI NGUIDJOÏ Thierry Emmanuel
% Version  : 1.0 -- Juin 2026
% Statut   : RÉDIGÉ
% Sources  : Ch. 9 (ADB, 4 niveaux acquisition NIST, iOS contraintes)
%            M5_Investigation_Numerique_PNC.docx (BLOC mobile)
%            Affaires MBONGO (téléphone EV-002), BAOBAB (MoMo fraud)
% Positionnement : Ch. 9 a couvert l'ACQUISITION mobile (ADB pull,
%                  sauvegarde, 4 niveaux). Ce chapitre couvre l'ANALYSE :
%                  que faire de ce qu'on a extrait ? Bases de données SQLite
%                  des applications, artefacts WhatsApp, artefacts iOS,
%                  géolocalisation, artefacts Mobile Money.
% ============================================================================

\chapter{Forensique Mobile --- Android et iOS}
\label{chap:for-mob}

\epigraph{%
    « Le téléphone portable est le journal intime du XXI\textsuperscript{e}
    siècle --- et bien plus révélateur. »%
}{--- Adage de la forensique mobile}

\niveauChapitre{Intermédiaire}{%
    Chapitre~\ref{chap:acq} (acquisition mobile~: ADB, 4~niveaux,
    iOS) indispensable. Ce chapitre couvre l'analyse post-acquisition.%
}

\begin{boxobjectifs}
\begin{itemize}
    \item Comprendre l'architecture de stockage Android~: partitions,
          systèmes de fichiers, emplacements des données applicatives.
    \item Analyser les bases de données SQLite des applications
          critiques~: SMS, WhatsApp, contacts, calendrier.
    \item Reconstruire une chronologie d'activité mobile depuis
          les journaux système Android et les métadonnées EXIF.
    \item Identifier et exploiter les artefacts de géolocalisation~:
          historique GPS, réseaux Wi-Fi, antennes GSM.
    \item Analyser les artefacts spécifiques Mobile Money (MTN MoMo,
          Orange Money) pour les fraudes financières numériques.
    \item Maîtriser les particularités forensiques d'iOS~: sauvegarde
          iTunes, données iCloud, limitations techniques.
\end{itemize}
\end{boxobjectifs}

\bigskip

% ============================================================================
\section*{Positionnement~: de l'extraction à l'analyse}
% ============================================================================

Le Chapitre~\ref{chap:acq} s'est terminé sur \texttt{adb pull /sdcard/}
et la liste des 9~commandes ADB essentielles. Ce chapitre part de là~:
vous avez l'extraction. Que contient-elle~? Où sont les preuves~? Comment
les lire, les dater, les corroborer~?

En Afrique subsaharienne, l'écrasante majorité des infractions numériques
implique un téléphone mobile plutôt qu'un ordinateur~: fraudes MoMo,
harcèlement sur WhatsApp, arnaques Facebook, SIM~swapping. Maîtriser
la forensique mobile n'est pas optionnel --- c'est la compétence la
plus immédiatement utile sur le terrain camerounais.

\separateur

% ============================================================================
\section{Architecture de Stockage Android}
\label{sec:architecture-android}
% ============================================================================

\subsection{Les partitions essentielles}

Un téléphone Android typique présente plusieurs partitions dont deux
sont centrales pour la forensique~:

\begin{tableaucours}{p{2.8cm}p{3.5cm}p{5.5cm}}{%
    \tableauHeader{Partition} &
    \tableauHeader{Contenu} &
    \tableauHeader{Accessibilité forensique}
}
\texttt{/data} & Données applicatives, comptes,
    bases de données SQLite, paramètres &
    \textbf{Nécessite root ou extraction physique.}
    Contient les données les plus sensibles~: messages,
    mots de passe, tokens d'authentification. \\
\texttt{/sdcard} (\texttt{/storage/emulated/0}) &
    Photos, vidéos, téléchargements, backups WhatsApp &
    \textbf{Accessible via ADB sans root.}
    \texttt{adb pull /sdcard/} \\
\texttt{/system} & OS Android, applications système &
    Lecture seule~; peu utile en investigation courante \\
\texttt{/cache} & Cache des applications &
    Volatile~; contenu non garanti après redémarrage \\
\end{tableaucours}
\vspace{-0.8em}
\begin{center}{\small\itshape Tableau~16.1 --- Partitions Android et accessibilité forensique}\end{center}

\begin{boxalert}
\textbf{La limite fondamentale de l'extraction ADB sans root}

\texttt{adb pull /sdcard/} donne accès à la carte SD virtuelle
(\texttt{/storage/emulated/0}). Les données applicatives les plus
importantes --- SMS, contacts, messages WhatsApp, tokens de connexion
--- sont dans \texttt{/data/data/<package>} et \textbf{ne sont pas
accessibles sans root}.

\textbf{Conséquence pratique~:} pour une affaire impliquant les messages
d'une application (WhatsApp, SMS, Messenger), l'extraction ADB logique
est insuffisante sauf si~: (a)~le téléphone était rooté, (b)~une
sauvegarde ADB a été activée~(\texttt{adb backup -all -shared}), ou
(c)~les backups WhatsApp ont été sauvegardés sur \texttt{/sdcard}
(comportement par défaut de WhatsApp).
\end{boxalert}

\subsection{Emplacements clés des données applicatives}

\begin{tableaucours}{p{3.5cm}p{4.0cm}p{4.0cm}}{%
    \tableauHeader{Application} &
    \tableauHeader{Chemin (root requis)} &
    \tableauHeader{Chemin accessible sans root}
}
SMS & \texttt{/data/data/com.android.providers.telephony/databases/mmssms.db} &
    \texttt{adb shell content query --uri content://sms/} \\
Contacts & \texttt{/data/data/com.android.providers.contacts/databases/contacts2.db} &
    \texttt{adb shell content query --uri content://contacts/} \\
WhatsApp messages & \texttt{/data/data/com.whatsapp/databases/msgstore.db} &
    \texttt{/sdcard/WhatsApp/Databases/msgstore.db.crypt15} (chiffré) \\
WhatsApp médias & \texttt{/data/data/com.whatsapp/} &
    \texttt{/sdcard/WhatsApp/Media/} (non chiffré) \\
Appels & \texttt{/data/data/com.android.providers.contacts/databases/contacts2.db} &
    \texttt{adb shell content query --uri content://call\_log/} \\
Localisation (Google) & \texttt{/data/data/com.google.android.gms/} &
    Google Timeline via réquisition (legalprocess.google.com) \\
\end{tableaucours}
\vspace{-0.8em}
\begin{center}{\small\itshape Tableau~16.2 --- Emplacements des données applicatives}\end{center}

% ============================================================================
\section{Analyse des Bases de Données SQLite}
\label{sec:sqlite}
% ============================================================================

Les applications Android stockent la quasi-totalité de leurs données
dans des fichiers SQLite (extension \texttt{.db}). Ces bases de données
sont le c\oe ur de l'analyse forensique mobile.

\subsection{Analyser les SMS~: \texttt{mmssms.db}}

\begin{boxoutils}{Analyse de mmssms.db --- base de données SMS}
\begin{lstlisting}[language=bash_forensique]
# Ouverture de la base de donnees SMS (extraite via ADB ou backup)
sqlite3 /output/extraction/mmssms.db

# Structure de la table SMS
.schema sms
-- CREATE TABLE sms (
--   _id INTEGER PRIMARY KEY,
--   thread_id INTEGER,  -- identifiant conversation
--   address TEXT,       -- numero de telephone
--   date INTEGER,       -- timestamp Unix en millisecondes
--   date_sent INTEGER,
--   protocol INTEGER,
--   read INTEGER,       -- 1=lu, 0=non lu
--   type INTEGER,       -- 1=recu, 2=envoye, 3=brouillon
--   body TEXT,          -- contenu du message
-- )

-- Tous les SMS avec horodatage lisible
SELECT
    datetime(date/1000, 'unixepoch', 'localtime') AS date_heure,
    address AS correspondant,
    CASE type
        WHEN 1 THEN 'Recu'
        WHEN 2 THEN 'Envoye'
        WHEN 3 THEN 'Brouillon'
    END AS type_sms,
    body AS contenu
FROM sms
ORDER BY date;

-- SMS contenant des mots-cles suspects
SELECT datetime(date/1000, 'unixepoch') AS date,
       address, body
FROM sms
WHERE body LIKE '%code%' OR body LIKE '%OTP%'
   OR body LIKE '%virement%' OR body LIKE '%MoMo%'
ORDER BY date;
\end{lstlisting}
\end{boxoutils}

\subsection{Analyser WhatsApp~: \texttt{msgstore.db}}

\begin{boiteRemarque}{Le déchiffrement des backups WhatsApp}
WhatsApp stocke ses backups locaux au format \texttt{msgstore.db.crypt15}
(chiffrement AES-256-GCM). La clé est stockée dans
\texttt{/data/data/com.whatsapp/files/key} --- inaccessible sans root.

\textbf{Cas où le backup est analysable sans root~:}
\begin{enumerate}
    \item Téléphone \textbf{rooté}~: extraire directement \texttt{msgstore.db}
          depuis \texttt{/data/data/com.whatsapp/databases/}.
    \item \textbf{Sauvegarde ADB} (si activée)~: \texttt{adb backup
          -apk com.whatsapp} extrait l'application et ses données.
    \item \textbf{WhatsApp Web actif}~: la session Web peut révéler
          le contenu en temps réel (pas forensique, mais utile en
          surveillance légale autorisée).
    \item \textbf{Déchiffrement avec la clé extraite}~: si root obtenu
          via exploit, extraire \texttt{key} puis déchiffrer avec
          l'outil \texttt{whatsapp-viewer}.
\end{enumerate}
\end{boiteRemarque}

\begin{boxoutils}{WhatsApp --- base non chiffrée (root) ou médias}
\begin{lstlisting}[language=bash_forensique]
-- Analyse de msgstore.db (WhatsApp, necesssite root ou backup)
sqlite3 /output/msgstore.db

-- Messages recus et envoyes avec horodatage
SELECT
    datetime(timestamp/1000, 'unixepoch', 'localtime') AS date_heure,
    jid AS correspondant,    -- format : 237XXXXXXXXX@s.whatsapp.net
    CASE from_me
        WHEN 0 THEN 'Recu'
        WHEN 1 THEN 'Envoye'
    END AS direction,
    data AS message_texte
FROM messages
WHERE data IS NOT NULL
ORDER BY timestamp;

-- Groupes WhatsApp et participants
SELECT
    g.subject AS nom_groupe,
    m.jid AS membre,
    datetime(m.timestamp/1000, 'unixepoch') AS date_ajout
FROM group_participants_history m
JOIN chat g ON g.jid = m.gjid
ORDER BY g.subject, m.timestamp;

-- Medias envoyes/recus (avec chemins de fichier)
SELECT
    datetime(timestamp/1000, 'unixepoch') AS date,
    jid, message_type, media_caption, local_path
FROM messages
WHERE message_type IN (1,2,3)  -- image, audio, video
ORDER BY timestamp;
\end{lstlisting}

\textbf{Médias WhatsApp accessibles sans root~:}
\begin{lstlisting}[language=bash_forensique]
# Les medias (photos, videos, documents) envoyes/recus par WhatsApp
# sont stockes en CLAIR sur /sdcard/WhatsApp/Media/
ls /mnt/extraction/WhatsApp/Media/
# WhatsApp Images/
# WhatsApp Video/
# WhatsApp Documents/
# WhatsApp Audio/

# Hash SHA-256 de tous les medias
find /mnt/extraction/WhatsApp/Media/ -type f \
    -exec sha256sum {} \; > /output/whatsapp_media_hashes.txt

# Verification sur VirusTotal (documents/APK suspects)
# Extraction des metadonnees EXIF des images
exiftool /mnt/extraction/WhatsApp/Media/WhatsApp\ Images/*.jpg \
    | grep -i "GPS\|Date\|Camera\|Make"
\end{lstlisting}
\end{boxoutils}

\subsection{Analyser les journaux d'appels}

\begin{lstlisting}[language=bash_forensique,
    caption={Extraction des journaux d'appels}]
# Via ADB content query (sans root)
adb shell content query --uri content://call_log/calls \
    --projection "number,type,date,duration,name" \
    > /output/call_log.txt

# Interpretation du champ type :
# 1 = Appel entrant
# 2 = Appel sortant
# 3 = Appel manque
# 4 = Messagerie vocale

# Analyse des appels vers un numero suspect
adb shell content query --uri content://call_log/calls \
    --where "number='23767XXXXXXX'" \
    --projection "date,type,duration"

# Convertir les timestamps
# date en millisecondes Unix -> datetime
python3 -c "
from datetime import datetime
with open('/output/call_log.txt') as f:
    for line in f:
        if 'date' in line:
            ts = int(line.split('=')[1].split(',')[0])
            print(datetime.fromtimestamp(ts/1000))
"
\end{lstlisting}

% ============================================================================
\section{Géolocalisation et Mobilité}
\label{sec:geolocalisation}
% ============================================================================

\subsection{Métadonnées EXIF des photos}

Les photos prises avec un smartphone embarquent souvent les coordonnées
GPS dans les métadonnées EXIF~--- une source précieuse pour établir
où se trouvait le suspect (ou la victime) à un moment donné.

\begin{boxoutils}{Extraction et visualisation des métadonnées EXIF}
\begin{lstlisting}[language=bash_forensique]
# Extraire les metadonnees GPS de toutes les photos
exiftool -csv -GPS* /mnt/extraction/DCIM/Camera/*.jpg \
    > /output/photo_gps.csv

# Contenu type du CSV :
# SourceFile,GPSLatitude,GPSLongitude,GPSAltitude,DateTimeOriginal
# IMG_20260313_021705.jpg,3.8480 N,11.5021 E,741 m,2026:03:13 02:17:05

# Convertir en format decimal pour Google Maps
exiftool -p '$GPSLatitude#,$GPSLongitude#,$DateTimeOriginal $FileName' \
    /mnt/extraction/DCIM/Camera/*.jpg

# Photo prise a 02:17:05 le 13 mars 2026 -> coordonnees GPS
# Utiliser Google Maps ou QGIS pour visualiser le trajet
\end{lstlisting}
\end{boxoutils}

\begin{boxscenario}{Cas MBONGO --- localisation via EXIF}
L'extraction ADB du téléphone de MBONGO Joseph (EV-002,
Chapitre~\ref{chap:custody}) révèle 15~photos dans
\texttt{/sdcard/WhatsApp/Media/WhatsApp Images/}. L'analyse EXIF
de \texttt{IMG-20260313-WA0003.jpg} (une photo d'un relevé bancaire)
montre~:

\begin{lstlisting}[language=bash_forensique]
exiftool IMG-20260313-WA0003.jpg
# GPS Latitude   : 3 deg 51' 37.50" N   (3.8604 N)
# GPS Longitude  : 11 deg 30' 08.40" E  (11.5023 E)
# Create Date    : 2026:01:15 14:28:33
# Camera Model   : Samsung Galaxy A53
\end{lstlisting}

\textbf{Constatation~:} la photo \texttt{IMG-20260313-WA0003.jpg}
a été prise le 15~janvier~2026 à 14h28~min~33s à des coordonnées GPS
de \texttt{3.8604°N / 11.5023°E}, correspondant au quartier Akwa,
Douala (vérification sur OpenStreetMap). Le suspect affirmait être à
Yaoundé ce jour-là. La date et les coordonnées sont en cohérence directe
avec l'horodatage de l'email de phishing analysé dans l'affaire WOURI
(14h28~min, même secteur géographique).
\end{boxscenario}

\subsection{Historique de localisation Google}

Google enregistre la position des appareils Android dans Google Maps
Timeline (si la fonction est activée). Ces données sont accessibles
de deux façons~:

\begin{enumerate}
    \item \textbf{Via le compte Google du suspect} (si ses identifiants
          sont connus ou si son téléphone est déverrouillé)~:
          \texttt{myactivity.google.com} ou \texttt{timeline.google.com}.
    \item \textbf{Via réquisition légale à Google}~:
          \texttt{legalprocess.google.com} --- délai~: 10--30~jours~;
          base légale~: art.~57 \loicam{2010/012} (réquisition aux
          fournisseurs de services numériques).
\end{enumerate}

\subsection{Données réseau~: Wi-Fi et antennes GSM}

\begin{boxoutils}{Artefacts réseau sur Android}
\begin{lstlisting}[language=bash_forensique]
# Historique des reseaux Wi-Fi connus (sans root)
adb shell cat /data/misc/wifi/WifiConfigStore.xml
# Ou accessible dans le backup ADB :
# /apps/com.android.providers.settings/db/settings.db

# Analyse du fichier wpa_supplicant (reseaux Wi-Fi)
# Contient : SSID, BSSID, mot de passe (hash), dates de connexion
# Chaque SSID peut etre localise via les bases BSSID publiques
# (Wigle.net, Google Maps Wi-Fi API)

# Journaux des antennes GSM (Cell ID log)
# Disponible via Google Location History (réquisition)
# Ou sur certains appareils dans /data/logger/
# Format : timestamp, MCC, MNC, LAC, Cell_ID
# Convertir Cell_ID en position via opencellid.org

# Exemple de requete OpenCelliD API
curl "https://opencellid.org/cell/get?key=YOUR_KEY
    &mcc=624&mnc=01&lac=9001&cellid=12345"
# -> latitude, longitude de l'antenne (precision : 50-1000 m)
\end{lstlisting}
\end{boxoutils}

\begin{boxjuris}
\textbf{Données de localisation et vie privée~: cadre légal camerounais}

Les données de localisation (GPS, antennes GSM, Wi-Fi) constituent des
données à caractère personnel au sens de l'art.~5 de la \loicam{2024/017}
(\textit{toute information permettant d'identifier directement ou
indirectement une personne physique}). Leur collecte et traitement
en contexte forensique doivent respecter~:
\begin{itemize}
    \item Art.~52 \loicam{2010/012}~: autorisation judiciaire préalable~;
    \item Art.~57 \loicam{2010/012}~: réquisition formelle à l'opérateur~;
    \item Art.~9--11 \loicam{2024/017}~: finalité et proportionnalité.
\end{itemize}
\textbf{Point de vigilance~:} l'historique de localisation obtenu depuis
le téléphone saisi est admissible (l'appareil a été légalement saisi).
Mais demander l'historique de localisation à l'opérateur ou à Google
nécessite une réquisition judiciaire distincte.
\end{boxjuris}

% ============================================================================
\section{Artefacts Mobile Money~: MoMo, Orange Money}
\label{sec:mobile-money}
% ============================================================================

Les fraudes Mobile Money représentent la majorité des affaires numériques
en Afrique subsaharienne. La forensique mobile appliquée à ces fraudes
présente des spécificités importantes.

\subsection{Où sont stockées les données Mobile Money}

\begin{tableaucours}{p{3.5cm}p{4.0cm}p{4.5cm}}{%
    \tableauHeader{Opérateur / Données} &
    \tableauHeader{Source primaire} &
    \tableauHeader{Source secondaire (réquisition)}
}
MTN MoMo --- SMS de confirmation & \texttt{mmssms.db} &
    MTN Cameroun~; art.~57 Loi~2010 \\
MTN MoMo --- Historique transactions &
    Application MTN MoMo (\texttt{/data/data/com.mtn.mobile.cm/}) &
    Réquisition MTN Cameroun (CELLCOM) \\
Orange Money --- SMS &
    \texttt{mmssms.db} (SMS OTP, confirmations) &
    Orange Cameroun~; art.~57 Loi~2010 \\
Orange Money --- Transactions &
    Application Orange Money &
    Réquisition Orange Cameroun \\
Historique USSD &
    Non stocké sur l'appareil & \textbf{Uniquement via l'opérateur}~;
    réquisition indispensable \\
\end{tableaucours}
\vspace{-0.8em}
\begin{center}{\small\itshape Tableau~16.3 --- Sources des données Mobile Money}\end{center}

\begin{boxoutils}{Extraction des SMS MoMo depuis mmssms.db}
\begin{lstlisting}[language=bash_forensique]
-- Tous les SMS de MTN MoMo (numeroMTN = +237 670 000 000)
SELECT
    datetime(date/1000, 'unixepoch', 'localtime') AS date_heure,
    address AS expediteur,
    body AS contenu_sms
FROM sms
WHERE address = '670000000'      -- numero MoMo MTN Cameroun
   OR address LIKE '+23767%'
   OR body LIKE '%MoMo%' OR body LIKE '%MTN%'
ORDER BY date;

-- Resultat type (fraude par phishing OTP) :
-- 2026-01-15 14:22:11  670000000
--   "Votre code OTP MoMo est 847293. Ne le partagez jamais."
-- 2026-01-15 14:22:45  670000000
--   "Vous avez envoye 250000 FCFA a 677XXXXXX. Solde : 12500 FCFA"

-- Reconstituer la chronologie d'une fraude MoMo
SELECT
    datetime(date/1000, 'unixepoch', 'localtime') AS date_heure,
    CASE type
        WHEN 1 THEN 'SMS recu'
        WHEN 2 THEN 'SMS envoye'
    END AS direction,
    body
FROM sms
WHERE (body LIKE '%OTP%' OR body LIKE '%code%'
    OR body LIKE '%envoye%' OR body LIKE '%recu%')
  AND date BETWEEN
    strftime('%s','2026-01-15 14:00:00')*1000 AND
    strftime('%s','2026-01-15 15:00:00')*1000
ORDER BY date;
\end{lstlisting}
\end{boxoutils}

\begin{boxPNC}
\textbf{Réquisition aux opérateurs Mobile Money~: que demander}

La réquisition judiciaire (art.~57 \loicam{2010/012}) à MTN ou Orange
pour une fraude MoMo doit précisément demander~:
\begin{enumerate}
    \item \textbf{Historique des transactions} du numéro victime
          sur la période (30~jours recommandés + 15~jours avant)~;
    \item \textbf{Identité et photo} associées au numéro bénéficiaire
          (numéro ayant reçu les fonds)~;
    \item \textbf{Logs de connexion} à l'application ou aux services
          USSD du numéro bénéficiaire~;
    \item \textbf{Carte SIM associée}~: série SIM (ICCID) et historique
          des échanges SIM (SIM~swapping)~;
    \item \textbf{IMEI} du téléphone utilisé avec ce numéro.
\end{enumerate}
\textbf{L'IMEI est la clé~:} si le fraudeur a changé de SIM mais
gardé son téléphone, l'IMEI identifie le téléphone physique à travers
les changements de numéro.
\end{boxPNC}

% ============================================================================
\section{Forensique iOS~: Spécificités et Limitations}
\label{sec:forensique-ios}
% ============================================================================

\subsection{Architecture et contraintes forensiques}

iOS présente une architecture de sécurité qui complique significativement
la forensique par rapport à Android~:

\begin{tableaucours}{p{3.5cm}p{4.0cm}p{4.0cm}}{%
    \tableauHeader{Mécanisme iOS} &
    \tableauHeader{Impact forensique} &
    \tableauHeader{Contournement}
}
Chiffrement hardware AES-256 & Toutes les données sont chiffrées
    sur la puce~; inaccessibles sans code PIN & Nécessite
    Cellebrite UFED ou GrayKey~; ou accès au code PIN \\
Secure Enclave & Stocke les clés de déchiffrement~;
    inaccessible même avec root &
    Impossible à contourner sans exploit spécifique \\
App Sandbox & Chaque application est isolée~;
    aucun accès inter-app sans jailbreak &
    Exploitation jailbreak (iOS 16--17) \\
Effacement après 10~erreurs & 10~tentatives PIN incorrectes
    $\to$ effacement total &
    \textbf{Ne jamais tenter le brute-force PIN manuellement} \\
\end{tableaucours}
\vspace{-0.8em}
\begin{center}{\small\itshape Tableau~16.4 --- Contraintes forensiques iOS}\end{center}

\subsection{Sauvegarde iTunes/Finder~: extraction logique iOS}

Si le téléphone est \textbf{déverrouillé} et que l'utilisateur fait
confiance à l'ordinateur de l'investigateur~:

\begin{boxoutils}{Extraction via sauvegarde iTunes (macOS/Windows)}
\begin{lstlisting}[language=bash_forensique]
# PREREQUIS : telephone deverrouille + "Faire confiance" accepte

# macOS -- sauvegarde via Finder
# Ouvrir Finder -> iPhone -> Sauvegarder maintenant
# Emplacement : ~/Library/Application Support/MobileSync/Backup/

# Windows -- sauvegarde via iTunes
# Emplacement : %APPDATA%\Apple Computer\MobileSync\Backup\

# Extraire la sauvegarde avec iMazing (ou outil forensique)
# ou parser manuellement avec libimobiledevice :
idevicebackup2 backup /output/ios_backup/

# Parser la sauvegarde iOS avec iBackupBot ou iPhone Backup Extractor
# Structure de la sauvegarde :
# - Manifest.db (SQLite) : inventaire de tous les fichiers
# - Manifest.plist : informations de l'appareil
# - <domain>/<hash> : fichiers individuels

# Interroger le Manifest.db
sqlite3 ~/Library/Application\ Support/MobileSync/Backup/<UDID>/Manifest.db
SELECT relativePath, domain, flags, fileID
FROM Files
WHERE relativePath LIKE '%Messages%'
ORDER BY relativePath;
\end{lstlisting}
\end{boxoutils}

\subsection{Analyse des bases iOS~: Messages et iCloud}

\begin{boxoutils}{Analyse de sms.db iOS (messages iMessage et SMS)}
\begin{lstlisting}[language=bash_forensique]
-- Localiser le fichier Messages dans la sauvegarde
-- Domain : HomeDomain, relativePath : Library/SMS/sms.db

-- Extraire les messages
sqlite3 /output/ios_backup/sms.db

SELECT
    datetime(date/1000000000 + 978307200, 'unixepoch',
        'localtime') AS date_heure,
    -- iOS date = secondes depuis 01/01/2001 (pas Unix epoch)
    CASE is_from_me
        WHEN 0 THEN h.id   -- expediteur (numero)
        WHEN 1 THEN 'Moi'
    END AS expediteur,
    text AS message
FROM message m
LEFT JOIN handle h ON h.rowid = m.handle_id
ORDER BY date;

-- Attention a l'epoch iOS : +978307200 secondes (offset 2001 vs 1970)

-- Contacts iOS (AddressBook.sqlitedb)
SELECT
    ABPerson.First, ABPerson.Last,
    ABMultiValue.value AS telephone
FROM ABPerson
JOIN ABMultiValue ON ABMultiValue.record_id = ABPerson.rowid
WHERE ABMultiValue.property = 3  -- propriete telephone
ORDER BY ABPerson.Last;
\end{lstlisting}
\end{boxoutils}

\begin{boiteRemarque}{L'epoch iOS~: piège fréquent}
iOS utilise une epoch différente de Unix~: les timestamps iOS sont
comptés \textbf{depuis le 1er~janvier~2001} (et non le
1er~janvier~1970). L'offset est de 978~307~200~secondes.
Ne pas appliquer cette conversion donne des dates situées en 1970 ---
erreur immédiatement identifiable mais qui fragilise un rapport non relu.
Pour les messages (iMessage), il faut en plus diviser par~$10^9$
(nanoseconde~$\to$ seconde).
\end{boiteRemarque}

% ============================================================================
\section{Autopsy pour la Forensique Mobile}
\label{sec:autopsy-mobile}
% ============================================================================

\termecle{Autopsy}\index{Autopsy} (Brian Carrier, Basis Technology)
intègre depuis la version~4 des modules d'analyse mobile qui automatisent
une grande partie des analyses SQLite présentées dans ce chapitre.

\begin{boxoutils}{Autopsy --- analyse d'une extraction Android}
\begin{lstlisting}[language=bash_forensique]
# Workflow Autopsy pour forensique mobile :
# 1. New Case -> Add Data Source -> Logical Files
#    -> Selectionner le repertoire d'extraction ADB

# 2. Modules a activer :
#    - Android Analyzer : parse automatiquement mmssms.db, call_log,
#      contacts, browser history
#    - EXIF Parser : extrait toutes les metadonnees GPS
#    - Keyword Search : recherche de termes (numeros, noms, montants)
#    - Communication : reconstruit les conversations

# 3. Resultats automatiques :
#    Reports -> Data Artifacts -> Text Messages
#                              -> Call Logs
#                              -> Contacts
#                              -> GPS Tracks (carte interactive)

# Limitation importante :
# Autopsy ne peut parser que ce qui est ACCESSIBLE
# (extraction /sdcard)
# Pour /data/data (root requis), il faut l'image complete du FS
# via extraction physique (Cellebrite, UFED)

# Export du rapport HTML pour le dossier judiciaire :
# Generate Report -> HTML Report
# Inclut horodatages, sources, et résumé par artefact
\end{lstlisting}
\end{boxoutils}

\separateur

\begin{boxresume}
\textbf{Ce qu'il faut retenir de ce chapitre~:}
\begin{itemize}
    \item \textbf{Limite ADB sans root~:} \texttt{adb pull /sdcard/}
          donne les médias et backups WhatsApp, mais pas les messages
          chiffrés ni les données \texttt{/data/data}. La sauvegarde
          \texttt{adb backup -all -apk} peut donner accès aux données
          applicatives selon les réglages de l'application.

    \item \textbf{SQLite = c\oe ur de l'analyse mobile~:} SMS
          (\texttt{mmssms.db}), WhatsApp (\texttt{msgstore.db}),
          appels, contacts, navigateur -- tout est en SQLite.
          Les requêtes de ce chapitre sont directement réutilisables
          en investigation.

    \item \textbf{WhatsApp~:} médias en clair sur \texttt{/sdcard/WhatsApp/Media/}~;
          messages chiffrés dans \texttt{.crypt15}. Clé dans
          \texttt{/data/data/com.whatsapp/files/key} (root requis).
          Documenter la situation dans le rapport~: ``messages non
          déchiffrés faute de clé accessible''.

    \item \textbf{EXIF GPS}~: chaque photo peut contenir des
          coordonnées GPS précises. Utiliser \texttt{exiftool -csv -GPS*}
          pour extraire en masse~; visualiser sur OpenStreetMap.

    \item \textbf{Mobile Money}~: fraudes MoMo~$\to$ toujours
          analyser \texttt{mmssms.db} (SMS OTP~; confirmations
          de transactions) ET requérir les logs opérateur
          (art.~57 Loi~2010). L'IMEI du téléphone fraudeur
          est la clé de l'identification même si la SIM a changé.

    \item \textbf{iOS~:} ne jamais tenter le brute-force PIN
          (10~erreurs~$\to$ effacement). Extraction via sauvegarde
          iTunes si téléphone déverrouillé. Attention à l'epoch
          iOS~: offset +978~307~200~secondes vs Unix.

    \item \textbf{Autopsy~:} automatise l'analyse des extractions
          Android (SMS, appels, GPS, contacts, navigateur). Limitation~:
          ne traite que ce qui est accessible (pas le \texttt{/data/data}
          sans root).
\end{itemize}
\end{boxresume}

\bigskip

\begin{boxexo}[title={\ding{45}\quad Exercice~16.1 --- Requêtes SQLite \niveaubase}]
Vous disposez de \texttt{mmssms.db} extrait du téléphone de BELLO Armand
(Opération TANGUE). Rédigez les requêtes SQLite permettant de~:
\begin{enumerate}[(a)]
    \item Lister tous les SMS reçus entre le 1\textsuperscript{er} et
          le 15~mars~2026, avec horodatage lisible et expéditeur.
    \item Extraire tous les SMS contenant les mots ``code'', ``OTP'',
          ou ``virement''.
    \item Calculer le nombre de SMS échangés avec le numéro
          \texttt{+237677XXXXXX} par jour (un résultat par jour).
    \item Identifier tous les SMS d'un expéditeur qui a envoyé
          plus de 50~messages en un seul jour (indicateur d'automatisation).
\end{enumerate}
\end{boxexo}

\begin{boxexo}[title={\ding{45}\quad Exercice~16.2 --- Fraude MoMo \niveaumoyen}]
Dans l'affaire BAOBAB, la victime FOUDA Clément a perdu 250~000~FCFA.
L'analyse de son téléphone révèle les SMS suivants dans \texttt{mmssms.db}~:
\begin{lstlisting}[language=bash_forensique]
-- Extrait :
date_heure            expediteur    contenu
2026-01-15 14:22:11   670000000
    "Votre code OTP MoMo est 847293. Ne le partagez jamais."
2026-01-15 14:22:15   +237677000123
    "Salut. Je travaille chez MTN. Donnez-moi votre code pour
     verifier votre compte."
2026-01-15 14:22:45   670000000
    "Vous avez envoye 250000 FCFA a 677000123. Solde : 3400 FCFA"
\end{lstlisting}
\begin{enumerate}[(a)]
    \item Qualifiez juridiquement les infractions commises (articles
          précis des Lois~2010 et~2024).
    \item Que demandez-vous dans la réquisition à MTN Cameroun~?
          Listez les 5~informations indispensables.
    \item Comment l'IMEI peut-il aider l'enquête si le fraudeur
          a changé de SIM après la fraude~?
    \item Rédigez la constatation forensique correspondante à partir
          des SMS ci-dessus, en appliquant la distinction
          constatation/interprétation.
\end{enumerate}
\end{boxexo}

\begin{boxexo}[title={\ding{45}\quad Exercice~16.3 --- Investigation mobile complète \niveauexpert}]
Vous avez en custody le téléphone Samsung Galaxy A53 de MBONGO Joseph
(EV-002, Chapitre~\ref{chap:custody}). L'extraction ADB a produit
les données suivantes dans \texttt{/output/extraction\_EV002/}~:
\begin{itemize}
    \item \texttt{/sdcard/} (photos, téléchargements, backups WhatsApp)~;
    \item \texttt{mmssms.db}, \texttt{call\_log.db}, \texttt{contacts2.db}~;
    \item \texttt{/sdcard/WhatsApp/Databases/msgstore.db.crypt15}~;
    \item \texttt{/sdcard/WhatsApp/Media/WhatsApp Images/} (23~images).
\end{itemize}
\begin{enumerate}[(a)]
    \item Planifiez l'analyse~: dans quel ordre analyser ces sources~?
          Quelle question investigative posez-vous à chaque source~?
    \item Les messages WhatsApp sont dans \texttt{.crypt15} (chiffrés).
          Quelles sont vos options légales et techniques pour y accéder~?
          Que documentez-vous si l'accès est impossible~?
    \item \texttt{exiftool} révèle que la photo
          \texttt{IMG-20260115-WA0003.jpg} a des coordonnées GPS
          situées à 200~mètres du siège de SONATRAV~SA à 14h28.
          L'affaire principale concerne un virement frauduleux
          depuis les comptes de SONATRAV. Quelle valeur probatoire
          accordez-vous à cette photo~? Rédigez la constatation et
          l'analyse correspondantes.
    \item Rédigez la section ``Artefacts mobiles'' du rapport d'expertise,
          couvrant les 4~sources analysées, avec pour chacune~:
          méthode d'extraction, artefacts trouvés, valeur probatoire.
\end{enumerate}
\end{boxexo}

\bigskip

\begin{boxinfo}
\textbf{Transition.} Le Chapitre~\ref{chap:anti-for} (Anti-Forensique
et Contre-Mesures) traite les techniques d'effacement et d'obfuscation
que les fraudeurs utilisent sur les téléphones~--- y compris le
factory reset, l'effacement des logs d'appels, les applications
``vault'' et les téléphones chiffrés. Ce chapitre vous apprendra
à détecter ces tentatives d'anti-forensique et à en documenter
l'existence même quand les données ont été effacées.
\end{boxinfo}
